% ============================================================================
% Tests report: exploratory link between video and pressure sensor.
% ============================================================================
\section{Contraste con el sensor de presión}\label{sec:presion}

Cada ensayo se registró también con el sensor de fuerza (FSR) del proyecto. Se usa su conductancia
$G = 1/R_{\mathrm{FSR}}$, que crece con la fuerza aplicada, y se la alinea con el video mediante el
pulso LED de sincronización (\cref{sec:escala_ens}). Este contraste es exploratorio: se compara la
\emph{movilidad} medida en el video (percentil 90 de la rapidez de los 22 robots) con la señal del
sensor, sin un modelo de cómo se transmite la fuerza hasta él.

\IfFileExists{generado/comparacion/presion.tex}{% Generated by comparar_ensayos.py
\begin{table}[H]
  \centering
  \caption{Alineación del registro de presión con el video y conductancia dentro y fuera de los atascos.}
  \label{tab:presion}
  \small
  \begin{tabular}{@{}lrrr@{}}
    \toprule
    & \Cmp{Lab}{A} & \Cmp{Lab}{B} & \Cmp{Lab}{C} \\
    \midrule
    Pulso LED: inicio [\si{\second}] & \num{\Cmp{Led}{A}} & \num{\Cmp{Led}{B}} & \num{\Cmp{Led}{C}} \\
    Pulso LED: fin ($=$ inicio del video) [\si{\second}] & \num{\Cmp{LedFin}{A}} & \num{\Cmp{LedFin}{B}} & \num{\Cmp{LedFin}{C}} \\
    Retardo por correlación cruzada [\si{\second}] & \num{\Cmp{Lag}{A}} & \num{\Cmp{Lag}{B}} & \num{\Cmp{Lag}{C}} \\
    Coeficiente de correlación $r$ & \num{\Cmp{Rxc}{A}} & \num{\Cmp{Rxc}{B}} & \num{\Cmp{Rxc}{C}} \\
    Correlación cruzada $-$ fin del pulso [\si{\second}] & \num{\Cmp{LagLed}{A}} & \num{\Cmp{LagLed}{B}} & \num{\Cmp{LagLed}{C}} \\
    $G$ mediana en atascos [\si{\micro\siemens}] & \num{\Cmp{Gin}{A}} & \num{\Cmp{Gin}{B}} & \num{\Cmp{Gin}{C}} \\
    $G$ mediana fuera de atascos [\si{\micro\siemens}] & \num{\Cmp{Gout}{A}} & \num{\Cmp{Gout}{B}} & \num{\Cmp{Gout}{C}} \\
    Cociente & \num{\Cmp{Grat}{A}} & \num{\Cmp{Grat}{B}} & \num{\Cmp{Grat}{C}} \\
    Muestras saturadas en atascos [\si{\percent}] & \num{\Cmp{SatIn}{A}} & \num{\Cmp{SatIn}{B}} & \num{\Cmp{SatIn}{C}} \\
    Muestras saturadas en total [\si{\percent}] & \num{\Cmp{SatAll}{A}} & \num{\Cmp{SatAll}{B}} & \num{\Cmp{SatAll}{C}} \\
    \bottomrule
  \end{tabular}
\end{table}
\begin{figure}[H]
  \centering
  \begin{tikzpicture}
    \begin{axis}[width=0.93\linewidth, height=3.6cm, axis y line*=left, ymin=0, ymax=8, xmin=0, enlarge x limits=false, ylabel={$|\vec v|_{90}$ [\si{\milli\metre\per\second}]}, grid=major, grid style={gray!20}, legend style={font=\footnotesize}, tick label style={font=\footnotesize}, title={\Cmp{Lab}{A}}, title style={font=\small, yshift=-1ex}]
      \addplot [ybar interval, fill=gray!30, draw=none] table [x=t, y expr=8*\thisrow{atasco}, col sep=comma] {generado/comparacion/presion_A.csv};
      \addplot [azul, thin] table [x=t, y=v_p90, col sep=comma] {generado/comparacion/presion_A.csv};
    \end{axis}
    \begin{axis}[width=0.93\linewidth, height=3.6cm, axis y line*=right, axis x line=none, xmin=0, enlarge x limits=false, ylabel={$\log_{10} G$ [\si{\micro\siemens}]}, ylabel style={rojo}, tick label style={font=\footnotesize}]
      \addplot [rojo, thin] table [x=t, y=logG, col sep=comma] {generado/comparacion/presion_A.csv};
    \end{axis}
  \end{tikzpicture}

  \begin{tikzpicture}
    \begin{axis}[width=0.93\linewidth, height=3.6cm, axis y line*=left, ymin=0, ymax=8, xmin=0, enlarge x limits=false, ylabel={$|\vec v|_{90}$ [\si{\milli\metre\per\second}]}, grid=major, grid style={gray!20}, legend style={font=\footnotesize}, tick label style={font=\footnotesize}, title={\Cmp{Lab}{B}}, title style={font=\small, yshift=-1ex}]
      \addplot [ybar interval, fill=gray!30, draw=none] table [x=t, y expr=8*\thisrow{atasco}, col sep=comma] {generado/comparacion/presion_B.csv};
      \addplot [azul, thin] table [x=t, y=v_p90, col sep=comma] {generado/comparacion/presion_B.csv};
    \end{axis}
    \begin{axis}[width=0.93\linewidth, height=3.6cm, axis y line*=right, axis x line=none, xmin=0, enlarge x limits=false, ylabel={$\log_{10} G$ [\si{\micro\siemens}]}, ylabel style={rojo}, tick label style={font=\footnotesize}]
      \addplot [rojo, thin] table [x=t, y=logG, col sep=comma] {generado/comparacion/presion_B.csv};
    \end{axis}
  \end{tikzpicture}

  \begin{tikzpicture}
    \begin{axis}[width=0.93\linewidth, height=3.6cm, axis y line*=left, ymin=0, ymax=8, xmin=0, enlarge x limits=false, ylabel={$|\vec v|_{90}$ [\si{\milli\metre\per\second}]}, grid=major, grid style={gray!20}, legend style={font=\footnotesize}, tick label style={font=\footnotesize}, title={\Cmp{Lab}{C}}, title style={font=\small, yshift=-1ex}, xlabel={$t$ de video [\si{\minute}]}]
      \addplot [ybar interval, fill=gray!30, draw=none] table [x=t, y expr=8*\thisrow{atasco}, col sep=comma] {generado/comparacion/presion_C.csv};
      \addplot [azul, thin] table [x=t, y=v_p90, col sep=comma] {generado/comparacion/presion_C.csv};
    \end{axis}
    \begin{axis}[width=0.93\linewidth, height=3.6cm, axis y line*=right, axis x line=none, xmin=0, enlarge x limits=false, ylabel={$\log_{10} G$ [\si{\micro\siemens}]}, ylabel style={rojo}, tick label style={font=\footnotesize}]
      \addplot [rojo, thin] table [x=t, y=logG, col sep=comma] {generado/comparacion/presion_C.csv};
    \end{axis}
  \end{tikzpicture}

  \caption{Movilidad de los robots (percentil 90 de la rapidez, azul; atascos sombreados) y conductancia del sensor de presión $G=1/R$ (rojo, escala logarítmica, medianas de \SI{10}{\second}), con el registro alineado al fin del pulso LED de sincronización (\cref{tab:presion}).}
  \label{fig:cmp_presion}
\end{figure}
}{}

\paragraph{Resultado.} En \Cmp{Lab}{A} la relación es directa: los dos atascos largos coinciden con
los intervalos de máxima conductancia (\cref{fig:cmp_presion}). Durante los atascos la mediana de $G$
es de \SI{\Cmp{Gin}{A}}{\micro\siemens}, contra \SI{\Cmp{Gout}{A}}{\micro\siemens} fuera de ellos, y el
\SI{\Cmp{SatIn}{A}}{\percent} de las muestras está saturado (el ADC llega a su máximo, $R_{\mathrm{FSR}}
= \SI{1141.5}{\ohm}$ con la resistencia de realimentación de \SI{5}{\kilo\ohm} que se usaba ese día). El
episodio de saturación de \SI{382}{\second} documentado en el análisis del sensor coincide con el
atasco de \SI{\Cmp{AtMax}{A}}{\second} entre los minutos 45 y 51 del video: la saturación empieza
unos \SI{13}{\second} después de que el grupo se detiene y termina junto con el atasco (dentro de él,
el \SI{96}{\percent} de las muestras está saturado). En los otros dos atascos de \Cmp{Lab}{A} el sensor
no satura, pero su conductancia mediana (\num{87} y \SI{250}{\micro\siemens}) sigue siendo miles de
veces mayor que fuera de los atascos. En \Cmp{Lab}{C} el atasco
principal (minutos 19 a 21) también coincide con un pico de conductancia, con $G$ \Cmp{Grat}{C} veces
mayor en los atascos que fuera de ellos. En \Cmp{Lab}{B}, en cambio, el único atasco no se ve en el
sensor (cociente \num{\Cmp{Grat}{B}}).

\paragraph{Interpretación.} Un atasco es un estado bloqueado: los robots siguen empujando, pero la red
de contactos no les deja lugar para moverse ni girar. Si el sensor queda dentro de esa red de
contactos, recibe una fuerza sostenida y alta; si el bloqueo se forma lejos del sensor, este no lo
registra, como en \Cmp{Lab}{B}. Los picos de $G$ fuera de los atascos (\cref{fig:cmp_presion})
corresponden a contactos transitorios que no llegan a detener a todo el grupo. Para afirmar más que
esta coincidencia haría falta conocer la posición del sensor en la arena y medir qué robots lo
tocan en cada momento, cosa que el video permite si se marca la ubicación del sensor.
